\documentclass{article}
\usepackage{graphicx} % Required for inserting images

\title{Speciale}
\author{xfh635 }
\date{March 2026}

\begin{document}

\maketitle

\section{Introduction}



\subsection{Category 5: Very High AI Dependency}

Category 5 contains firms for which AI was already economically central to the business model by 31 December 2022. Firms in this category are not merely active in AI-related markets or investing in AI capabilities; rather, AI was already a core driver of revenues, demand, monetization, or near-term growth expectations before the generative-AI boom of 2023. The classification is intentionally conservative. If the evidence suggested that AI was important but not yet clearly central to the firm’s economic model, the firm was assigned to a lower category.

Applying this strict ex ante criterion results in a very small set of firms. Only three S\&P 500 companies were judged to fit Category 5 with sufficiently high confidence: NVIDIA, Alphabet, and Meta Platforms. This narrow selection is methodologically desirable, as it creates a high-purity top category that can serve as a clear benchmark against firms with weaker AI dependency.

\begin{table}[H]
\centering
\caption{Category 5 firms: Very high AI dependency}
\begin{tabular}{|c|l|c|l|c|}
\hline
\textbf{Rank} & \textbf{Company} & \textbf{Ticker} & \textbf{Sector} & \textbf{Confidence score} \\ \hline
1 & NVIDIA & NVDA & Information Technology & 96 \\ \hline
2 & Alphabet & GOOGL & Communication Services & 88 \\ \hline
3 & Meta Platforms & META & Communication Services & 86 \\ \hline
\end{tabular}
\label{tab:category5}
\end{table}

\subsubsection*{NVIDIA}
\textbf{Ticker:} NVDA \\
\textbf{Sector:} Information Technology

NVIDIA is the clearest Category 5 firm in the S\&P 500. By the end of 2022, AI-related demand was already central to the company’s business, especially through its data center activities. NVIDIA’s product portfolio was directly tied to AI training and high-performance computing, and the company’s own reporting already identified artificial intelligence as a key part of its data center and compute business. This is important because it shows that NVIDIA’s AI exposure was not a speculative future opportunity but an economically material source of revenue before the 2023 boom. Relative to other semiconductor firms, NVIDIA had the most direct and concentrated exposure to AI infrastructure demand.

\textbf{Assessment:} Direct AI revenue exposure: High; Indirect AI demand exposure: High; Strategic importance of AI: High; Operational reliance on AI: Medium; Economic materiality of AI: High.

\textbf{Borderline note:} NVIDIA belongs in Category 5 rather than Category 4 because AI-related infrastructure demand was already central to the firm’s revenue base and growth profile by 2022, not merely an important adjacent opportunity.

\subsubsection*{Alphabet}
\textbf{Ticker:} GOOGL \\
\textbf{Sector:} Communication Services

Alphabet also fits Category 5 because AI was already deeply embedded in the company’s core monetization engine by the end of 2022. Search quality, ad targeting, ranking, recommendations, and automation were heavily dependent on machine learning and AI-based systems. In addition, AI was strategically important across Google Cloud and other products, but the key reason for the classification is that Alphabet’s core advertising business already relied on AI in an economically meaningful way. This distinguishes Alphabet from firms where AI was mainly a future growth lever or a product enhancement. By 2022, AI was already central to how Alphabet generated and optimized revenue at scale.

\textbf{Assessment:} Direct AI revenue exposure: Medium; Indirect AI demand exposure: High; Strategic importance of AI: High; Operational reliance on AI: High; Economic materiality of AI: High.

\textbf{Borderline note:} Alphabet is placed in Category 5 rather than Category 4 because AI was already a core part of the monetization structure of Search and digital advertising, rather than simply a major innovation priority.

\subsubsection*{Meta Platforms}
\textbf{Ticker:} META \\
\textbf{Sector:} Communication Services

Meta Platforms belongs in Category 5 because its advertising model already depended heavily on AI-driven ranking, targeting, recommendation, and measurement systems by the end of 2022. The firm’s revenues were overwhelmingly generated by digital advertising, and the effectiveness of this business model depended on its ability to match users with relevant content and advertisements. This made AI economically central to revenue generation rather than a secondary support function. Meta’s case differs from firms that used AI mainly for internal efficiency or product improvement. In Meta’s case, AI was already closely tied to the performance of the company’s dominant source of cash flow.

\textbf{Assessment:} Direct AI revenue exposure: Low; Indirect AI demand exposure: High; Strategic importance of AI: High; Operational reliance on AI: High; Economic materiality of AI: High.

\textbf{Borderline note:} Meta is placed in Category 5 rather than Category 4 because AI was already fundamental to the monetization and functioning of its core advertising platform, not merely a significant strategic capability.

\subsubsection*{Conclusion on Category 5}

These firms can be used as anchor firms for the highest AI-dependency category in the empirical analysis. Importantly, they were selected using only information available by 31 December 2022 and not based on subsequent stock price performance. The purpose of this category is therefore not to identify the firms that later benefited most from the generative-AI boom, but to isolate the firms for which AI was already economically central before that boom occurred. This makes Category 5 a useful benchmark for testing whether firms with the strongest ex ante AI dependency experienced larger stock price increases during the later period.

\subsection{Category 4: High AI Dependency}

Category 4 contains firms for which AI was already a major growth driver by 31 December 2022, but not the entire foundation of the business. These are companies where AI was materially shaping demand, product differentiation, commercialization, or customer adoption, while the firm’s total revenue base still depended on broader platforms, ecosystems, or adjacent business lines.

The classification is intentionally conservative. Firms were included only when the available pre-2023 evidence suggested that AI was already economically meaningful in a substantial way. If AI appeared to be strategically relevant but not yet clearly a major commercial driver, the firm was assigned to a lower category. Applying this standard results in a relatively small set of firms. Seven S\&P 500 companies were judged to fit Category 4 with sufficiently high confidence: Microsoft, Amazon, AMD, Adobe, Intuit, ServiceNow, and Arista Networks.

\begin{table}[H]
\centering
\caption{Category 4 firms: High AI dependency}
\begin{tabular}{|c|l|c|l|c|}
\hline
\textbf{Rank} & \textbf{Company} & \textbf{Ticker} & \textbf{Sector} & \textbf{Confidence score} \\ \hline
1 & Microsoft & MSFT & Information Technology & 90 \\ \hline
2 & Amazon & AMZN & Consumer Discretionary & 88 \\ \hline
3 & AMD & AMD & Information Technology & 86 \\ \hline
4 & Adobe & ADBE & Information Technology & 84 \\ \hline
5 & Intuit & INTU & Information Technology & 81 \\ \hline
6 & ServiceNow & NOW & Information Technology & 79 \\ \hline
7 & Arista Networks & ANET & Information Technology & 74 \\ \hline
\end{tabular}
\label{tab:category4}
\end{table}

\subsubsection*{Microsoft}
\textbf{Ticker:} MSFT \\
\textbf{Sector:} Information Technology

Microsoft fits Category 4 because, by the end of 2022, AI was already a major commercial driver within Azure and the broader Microsoft Cloud ecosystem. The company had integrated AI into cloud offerings and enterprise solutions in ways that materially supported growth, product differentiation, and customer adoption. At the same time, Microsoft remained a highly diversified company with large business lines in Office, Windows, gaming, security, and enterprise software. AI was therefore commercially important, but not yet the central economic foundation of the firm as a whole. This makes Microsoft a strong example of high AI dependency without meeting the stricter threshold for Category 5.

\textbf{Assessment:} Direct AI revenue exposure: Medium; Indirect AI demand exposure: High; Strategic importance of AI: High; Operational reliance on AI: High; Economic materiality of AI: High.

\textbf{Borderline note:} Microsoft belongs in Category 4 rather than Category 5 because AI was already a major growth driver, especially in cloud services, but was not yet clearly the dominant foundation of total company revenues and demand.

\textbf{Confidence score:} 90

\subsubsection*{Amazon}
\textbf{Ticker:} AMZN \\
\textbf{Sector:} Consumer Discretionary

Amazon belongs in Category 4 primarily because of AWS. By the end of 2022, AWS had already become a large and highly profitable business, and its offerings included machine learning and other tools that made AI-related demand commercially meaningful. This gave Amazon significant indirect exposure to AI adoption through the infrastructure layer of cloud computing. However, Amazon remained a diversified firm with major retail, logistics, subscription, and marketplace activities. AI was an important growth driver within AWS, but it was not the defining foundation of the consolidated business model. For that reason, Amazon is best treated as high AI dependency rather than very high AI dependency.

\textbf{Assessment:} Direct AI revenue exposure: Medium; Indirect AI demand exposure: High; Strategic importance of AI: High; Operational reliance on AI: Medium; Economic materiality of AI: High.

\textbf{Borderline note:} Amazon belongs in Category 4 rather than Category 5 because AI was economically meaningful through AWS, but the firm as a whole was still broader than AI-centric demand.

\textbf{Confidence score:} 88

\subsubsection*{AMD}
\textbf{Ticker:} AMD \\
\textbf{Sector:} Information Technology

AMD fits Category 4 because it had meaningful AI infrastructure exposure before the generative-AI boom, particularly through data center products and accelerators used for AI and machine learning workloads. The firm benefited from rising demand for compute capacity relevant to AI, yet its economic story in 2022 still extended well beyond AI alone. AMD’s growth also depended on broader server, processor, and semiconductor markets, which makes its exposure less concentrated than NVIDIA’s. In other words, AI was already commercially important and strategically relevant, but not yet the dominant economic driver of the company. This places AMD naturally in Category 4 rather than in the top category.

\textbf{Assessment:} Direct AI revenue exposure: Medium; Indirect AI demand exposure: High; Strategic importance of AI: High; Operational reliance on AI: Medium; Economic materiality of AI: High.

\textbf{Borderline note:} AMD belongs in Category 4 rather than Category 5 because AI demand was already important, but the firm’s 2022 revenue base still reflected broader semiconductor and server markets rather than clearly AI-centered dependence.

\textbf{Confidence score:} 86

\subsubsection*{Adobe}
\textbf{Ticker:} ADBE \\
\textbf{Sector:} Information Technology

Adobe belongs in Category 4 because AI was already embedded in the value proposition of its major software businesses by the end of 2022. Through Adobe Sensei, the company used AI to enhance creative workflows, automate document processes, and improve personalization in customer experience products. These capabilities were commercially meaningful because they supported product differentiation, user productivity, and enterprise adoption across Adobe’s core subscription software portfolio. However, Adobe’s revenues were still fundamentally generated by creative, document, and digital experience software rather than by AI as a stand-alone source of demand. AI was therefore a major growth and innovation driver, but not the central foundation of the company’s economics.

\textbf{Assessment:} Direct AI revenue exposure: Low; Indirect AI demand exposure: Medium; Strategic importance of AI: High; Operational reliance on AI: High; Economic materiality of AI: High.

\textbf{Borderline note:} Adobe belongs in Category 4 rather than Category 3 because AI was already embedded in large monetized product lines. It remains below Category 5 because the business was still fundamentally broader software rather than AI-centered demand.

\textbf{Confidence score:} 84

\subsubsection*{Intuit}
\textbf{Ticker:} INTU \\
\textbf{Sector:} Information Technology

Intuit is a strong Category 4 firm because, by 2022, AI was already closely tied to the company’s platform strategy and product economics. Across TurboTax, QuickBooks, and Credit Karma, the company used AI-driven prediction, automation, and personalization to improve user experience and support scalable service delivery. AI was not merely a supporting feature; it had become part of how Intuit positioned its platform and generated value for customers. At the same time, the company’s revenues still came from broader financial software and services rather than from AI itself. This makes AI highly important to Intuit’s growth profile and operating model, but not sufficiently dominant to justify Category 5.

\textbf{Assessment:} Direct AI revenue exposure: Low; Indirect AI demand exposure: Medium; Strategic importance of AI: High; Operational reliance on AI: High; Economic materiality of AI: High.

\textbf{Borderline note:} Intuit belongs in Category 4 rather than Category 3 because AI was already foundational to the firm’s platform strategy and commercial model. It remains below Category 5 because AI was not yet the central independent source of economic demand.

\textbf{Confidence score:} 81

\subsubsection*{ServiceNow}
\textbf{Ticker:} NOW \\
\textbf{Sector:} Information Technology

ServiceNow fits Category 4 because AI was already integrated into the commercial logic of its workflow automation platform by the end of 2022. The company used AI, machine learning, and predictive tools to enhance service management, workflow optimization, and enterprise productivity. These capabilities were economically meaningful because they improved the attractiveness of ServiceNow’s subscription offerings and reinforced the company’s core value proposition to enterprise customers. Nevertheless, the business was still fundamentally a digital workflow and enterprise software company rather than an AI-centered firm. AI was therefore a major driver of product quality and commercial expansion, but not the sole or dominant foundation of the firm’s revenue model.

\textbf{Assessment:} Direct AI revenue exposure: Low; Indirect AI demand exposure: Medium; Strategic importance of AI: High; Operational reliance on AI: High; Economic materiality of AI: High.

\textbf{Borderline note:} ServiceNow belongs in Category 4 rather than Category 3 because AI was already embedded in monetized flagship products and platform capabilities. It remains below Category 5 because the company’s economics still rested on a broader enterprise workflow subscription model.

\textbf{Confidence score:} 79

\subsubsection*{Arista Networks}
\textbf{Ticker:} ANET \\
\textbf{Sector:} Information Technology

Arista Networks is the weakest but still defensible inclusion in Category 4. By the end of 2022, the firm was benefiting from strong demand for high-performance networking infrastructure used in large-scale data centers and hyperscale environments, including workloads related to AI. This gives Arista meaningful indirect exposure to AI-related demand through the networking layer required for expanding compute intensity. However, the company remained fundamentally a networking infrastructure provider rather than a firm whose revenues were directly driven by AI products or AI-based monetization. For that reason, it sits near the lower boundary of Category 4, but still above Category 3 because AI-related infrastructure demand was already commercially relevant.

\textbf{Assessment:} Direct AI revenue exposure: Low; Indirect AI demand exposure: High; Strategic importance of AI: Medium; Operational reliance on AI: Medium; Economic materiality of AI: Medium.

\textbf{Borderline note:} Arista belongs in Category 4 rather than Category 3 because AI-related hyperscale networking demand was already commercially meaningful. It remains well below Category 5 because the company was still primarily a cloud-networking business rather than an AI-dependent business.

\textbf{Confidence score:} 74

\subsubsection*{Conclusion on Category 4}

These firms can be used as anchor firms for the high AI-dependency category in the empirical analysis. Importantly, they were selected using only information available by 31 December 2022 and not on the basis of subsequent stock price performance. The purpose of Category 4 is therefore to capture firms for which AI was already a major commercial driver before the generative-AI boom, but where the business model was still broader than AI alone. This makes the category a useful bridge between the very highest dependency firms in Category 5 and the larger set of firms where AI was present but less economically central.

\subsection{Category 3: Moderate AI Dependency}

Category 3 contains firms for which AI was already strategically relevant and economically meaningful by 31 December 2022, but not a dominant driver of the business. In this category, AI typically mattered through product enhancement, monetized features, workflow automation, customer value creation, or support for demand, while the firm’s revenues and near-term growth outlook still depended more on broader software, services, semiconductor, or cybersecurity franchises than on AI itself.

The classification is intentionally conservative. Firms were included only when the available pre-2023 evidence suggested that AI already had meaningful commercial relevance. If AI appeared to be a major growth driver, the firm was assigned to Category 4. If AI appeared only as a secondary feature or efficiency tool, the firm was assigned to Category 2. Applying this standard results in nine S\&P 500 firms being placed in Category 3: Salesforce, Oracle, Accenture, Qualcomm, Palo Alto Networks, Synopsys, Workday, Autodesk, and Fortinet.

\begin{table}[H]
\centering
\caption{Category 3 firms: Moderate AI dependency}
\begin{tabular}{|c|l|c|l|c|}
\hline
\textbf{Rank} & \textbf{Company} & \textbf{Ticker} & \textbf{Sector} & \textbf{Confidence score} \\ \hline
1 & Salesforce & CRM & Information Technology & 85 \\ \hline
2 & Oracle & ORCL & Information Technology & 83 \\ \hline
3 & Accenture & ACN & Information Technology & 81 \\ \hline
4 & Qualcomm & QCOM & Information Technology & 80 \\ \hline
5 & Palo Alto Networks & PANW & Information Technology & 79 \\ \hline
6 & Synopsys & SNPS & Information Technology & 78 \\ \hline
7 & Workday & WDAY & Information Technology & 76 \\ \hline
8 & Autodesk & ADSK & Information Technology & 74 \\ \hline
9 & Fortinet & FTNT & Information Technology & 73 \\ \hline
\end{tabular}
\label{tab:category3}
\end{table}

\subsubsection*{Salesforce}
\textbf{Ticker:} CRM \\
\textbf{Sector:} Information Technology

Salesforce fits Category 3 because, by the end of 2022, AI was already an important part of its CRM platform, but not the dominant economic foundation of the company. AI contributed to product functionality, recommendations, automation, and customer-facing enhancements across the broader cloud software portfolio. This made AI strategically relevant and commercially meaningful, but still as one component within a much larger subscription-based enterprise software model. Salesforce’s revenues continued to depend primarily on the breadth of its CRM and enterprise cloud offerings rather than on AI-specific demand. AI therefore had meaningful economic relevance, but not enough to classify the firm in Category 4.

\textbf{Assessment:} Direct AI revenue exposure: Low; Indirect AI demand exposure: Medium; Strategic importance of AI: High; Operational reliance on AI: Medium; Economic materiality of AI: Medium.

\textbf{Borderline note:} Salesforce belongs in Category 3 rather than Category 4 because AI was embedded in important customer-facing product capabilities, but it was not yet a major independent growth engine. It sits above Category 2 because AI was already commercially relevant, not merely an internal efficiency tool.

\textbf{Confidence score:} 85

\subsubsection*{Oracle}
\textbf{Ticker:} ORCL \\
\textbf{Sector:} Information Technology

Oracle belongs in Category 3 because AI and machine learning were already integrated into cloud applications and enterprise software offerings by the end of 2022, but the company remained economically much broader than AI. AI contributed to the value proposition of Oracle’s cloud and software stack, especially in analytics, enterprise applications, and automation. However, the company’s revenues still depended primarily on a wide range of database, infrastructure, and enterprise software activities rather than on AI itself. This makes AI strategically important and economically meaningful, but not dominant. Oracle therefore fits best in the moderate AI dependency category rather than in a higher bucket.

\textbf{Assessment:} Direct AI revenue exposure: Low; Indirect AI demand exposure: Medium; Strategic importance of AI: High; Operational reliance on AI: Medium; Economic materiality of AI: Medium.

\textbf{Borderline note:} Oracle belongs in Category 3 rather than Category 4 because AI was integrated into monetized offerings, but there is insufficient evidence that it was already a major separate growth driver by end-2022. It is above Category 2 because AI was already part of externally sold enterprise functionality.

\textbf{Confidence score:} 83

\subsubsection*{Accenture}
\textbf{Ticker:} ACN \\
\textbf{Sector:} Information Technology

Accenture fits Category 3 because AI was already commercially relevant to its consulting and digital transformation business by 2022, but not dominant enough to define the firm’s economics. AI mattered through client demand for automation, analytics, and digital transformation services, and it contributed to the company’s positioning in large enterprise projects. However, Accenture remained a broad professional-services company serving many industries and functions, with AI acting as one important capability among several. This makes AI strategically meaningful and economically relevant, but not strong enough to be considered a major stand-alone growth engine. Accenture therefore belongs in Category 3 rather than Category 4.

\textbf{Assessment:} Direct AI revenue exposure: Low; Indirect AI demand exposure: Medium; Strategic importance of AI: High; Operational reliance on AI: Medium; Economic materiality of AI: Medium.

\textbf{Borderline note:} Accenture belongs in Category 3 rather than Category 4 because AI supported service demand and differentiation, but did not yet make the firm heavily economically dependent on AI. It remains above Category 2 because AI was already part of client-facing revenue-generating work.

\textbf{Confidence score:} 81

\subsubsection*{Qualcomm}
\textbf{Ticker:} QCOM \\
\textbf{Sector:} Information Technology

Qualcomm belongs in Category 3 because AI was already relevant to its chip roadmap and product positioning in smartphones, IoT, automotive, and edge computing by the end of 2022. AI capabilities helped strengthen device performance and supported commercial relevance in several end markets. However, Qualcomm’s overall economics still depended primarily on its broader semiconductor and licensing model rather than on AI-specific demand. AI was therefore meaningful to product differentiation and industry positioning, but not the dominant driver of the company’s revenues or demand. This places Qualcomm naturally in the moderate AI dependency category rather than in a higher category.

\textbf{Assessment:} Direct AI revenue exposure: Low; Indirect AI demand exposure: Medium; Strategic importance of AI: High; Operational reliance on AI: Medium; Economic materiality of AI: Medium.

\textbf{Borderline note:} Qualcomm belongs in Category 3 rather than Category 4 because AI strengthened product relevance and future demand, but the firm remained primarily a connectivity, mobile chip, and licensing business. It remains above Category 2 because AI was already part of monetized product capability.

\textbf{Confidence score:} 80

\subsubsection*{Palo Alto Networks}
\textbf{Ticker:} PANW \\
\textbf{Sector:} Information Technology

Palo Alto Networks fits Category 3 because AI and machine learning were already embedded in core cybersecurity offerings by the end of 2022, but the company’s business model remained fundamentally broader than AI. AI contributed to threat detection, automation, analytics, and the effectiveness of multiple subscription-based security products. This made AI economically relevant and strategically important, especially because product quality in cybersecurity depends heavily on intelligent detection and response systems. Even so, the company remained primarily a cybersecurity platform provider rather than a firm whose revenues were driven mainly by AI demand. For that reason, the firm is best placed in Category 3 rather than in Category 4.

\textbf{Assessment:} Direct AI revenue exposure: Low; Indirect AI demand exposure: Medium; Strategic importance of AI: High; Operational reliance on AI: High; Economic materiality of AI: Medium.

\textbf{Borderline note:} Palo Alto Networks belongs in Category 3 rather than Category 4 because AI materially strengthened the value proposition of its security products, but was not yet a major independent growth engine. It is above Category 2 because AI was already central to product functionality.

\textbf{Confidence score:} 79

\subsubsection*{Synopsys}
\textbf{Ticker:} SNPS \\
\textbf{Sector:} Information Technology

Synopsys fits Category 3 because AI was meaningful on both the product side and the demand side by the end of 2022, but still not dominant enough to define the business. AI and machine learning mattered because they increased activity in semiconductor and system design, while also enhancing Synopsys’s own design automation tools. This gave the company both indirect exposure to AI-related chip demand and direct product relevance through AI-enabled design solutions. However, Synopsys remained fundamentally an electronic design automation and semiconductor IP company, with AI acting as an important commercial layer rather than the central economic foundation. This makes Category 3 the most appropriate classification.

\textbf{Assessment:} Direct AI revenue exposure: Low; Indirect AI demand exposure: Medium; Strategic importance of AI: High; Operational reliance on AI: Medium; Economic materiality of AI: Medium.

\textbf{Borderline note:} Synopsys belongs in Category 3 rather than Category 4 because AI was commercially relevant to both demand and product differentiation, but not yet strong enough to constitute a major growth engine. It is above Category 2 because AI already had visible commercial importance.

\textbf{Confidence score:} 78

\subsubsection*{Workday}
\textbf{Ticker:} WDAY \\
\textbf{Sector:} Information Technology

Workday belongs in Category 3 because AI and machine learning were already relevant to its cloud applications and customer value proposition by the end of 2022, but not dominant enough to determine the firm’s economics. AI helped improve analytics, automation, skills intelligence, and workflow efficiency in human capital management and financial software. These capabilities were commercially meaningful and strategically important, especially in enhancing the attractiveness of Workday’s subscription products. However, the company remained fundamentally dependent on its broader HCM and finance cloud application business rather than on AI-specific demand. This makes Workday a clear moderate AI dependency firm rather than a high AI dependency firm.

\textbf{Assessment:} Direct AI revenue exposure: Low; Indirect AI demand exposure: Medium; Strategic importance of AI: High; Operational reliance on AI: High; Economic materiality of AI: Medium.

\textbf{Borderline note:} Workday belongs in Category 3 rather than Category 4 because AI improved product capability and strategic direction, but the company still depended mainly on broader cloud software subscriptions. It remains above Category 2 because AI was already customer-facing and monetized indirectly.

\textbf{Confidence score:} 76

\subsubsection*{Autodesk}
\textbf{Ticker:} ADSK \\
\textbf{Sector:} Information Technology

Autodesk fits Category 3 because AI, machine learning, and generative design were already meaningful to its innovation strategy and product development by the end of 2022, but not yet dominant enough to define the company’s economics. AI contributed to automation, design insights, and workflow improvements across architecture, engineering, construction, and manufacturing tools. This made AI strategically relevant and commercially useful, especially in enhancing externally sold software products. However, Autodesk remained fundamentally a design software company, with revenues still driven by a wider portfolio of subscription products rather than by AI-specific demand. Category 3 is therefore more appropriate than either a lower or higher classification.

\textbf{Assessment:} Direct AI revenue exposure: Low; Indirect AI demand exposure: Low; Strategic importance of AI: High; Operational reliance on AI: Medium; Economic materiality of AI: Medium.

\textbf{Borderline note:} Autodesk belongs in Category 3 rather than Category 4 because AI was a meaningful innovation and product layer, but not a major stand-alone growth engine by end-2022. It is above Category 2 because AI was already embedded in customer-facing tools.

\textbf{Confidence score:} 74

\subsubsection*{Fortinet}
\textbf{Ticker:} FTNT \\
\textbf{Sector:} Information Technology

Fortinet is a lower-confidence but still defensible Category 3 firm. By the end of 2022, AI was already relevant to the company’s cybersecurity offerings, especially through subscription services and security operations solutions. AI contributed to threat detection, response, and the performance of commercial security services, making it economically meaningful rather than purely internal. At the same time, Fortinet remained fundamentally a broader network security and cybersecurity platform company, with revenues driven by hardware, subscriptions, and services across a wider security stack. AI was therefore strategically important and commercially useful, but not sufficiently dominant to qualify as a major growth engine. This places the firm in Category 3.

\textbf{Assessment:} Direct AI revenue exposure: Low; Indirect AI demand exposure: Medium; Strategic importance of AI: High; Operational reliance on AI: High; Economic materiality of AI: Medium.

\textbf{Borderline note:} Fortinet belongs in Category 3 rather than Category 4 because AI strengthened revenue-generating security offerings, but the evidence does not show AI as a major independent growth driver by end-2022. It is above Category 2 because AI was already tied to commercial products and services.

\textbf{Confidence score:} 73

\subsubsection*{Conclusion on Category 3}

These firms can be used as anchor firms for the moderate AI-dependency category in the empirical analysis. Importantly, they were selected using only information available by 31 December 2022 and not on the basis of later stock price performance. The purpose of Category 3 is therefore to capture firms for which AI was already visible in products, commercialization, customer value creation, or operating logic, but not yet dominant enough to define the business model. This makes the category useful as a middle group between the high AI-dependency firms in Categories 4 and 5 and the lower AI-dependency firms in Categories 1 and 2.

\subsection{Category 2: Low AI Dependency}

Category 2 contains firms for which AI was clearly present by 31 December 2022, but mainly as a secondary tool, product feature, analytics layer, or efficiency enhancer rather than a major driver of revenue, demand, or near-term growth. In this category, AI typically appears in fraud detection, personalization, customer service, workflow automation, logistics, or decision support, while the core business model remains largely unchanged without it.

The classification is intentionally conservative. Firms were included only when the available pre-2023 evidence suggested that AI was already visible in products, services, or operations, but not economically central enough to justify Category 3. If AI appeared strategically relevant and economically meaningful, the firm was assigned to Category 3. Applying this standard results in nine S\&P 500 firms being placed in Category 2: Mastercard, Bank of America, American Express, Starbucks, McDonald's, CVS Health, Lowe's, ADP, and UPS.

\begin{table}[H]
\centering
\caption{Category 2 firms: Low AI dependency}
\begin{tabular}{|c|l|c|l|c|}
\hline
\textbf{Rank} & \textbf{Company} & \textbf{Ticker} & \textbf{Sector} & \textbf{Confidence score} \\ \hline
1 & Mastercard & MA & Financials & 84 \\ \hline
2 & Bank of America & BAC & Financials & 82 \\ \hline
3 & American Express & AXP & Financials & 80 \\ \hline
4 & Starbucks & SBUX & Consumer Discretionary & 78 \\ \hline
5 & McDonald's & MCD & Consumer Discretionary & 77 \\ \hline
6 & CVS Health & CVS & Health Care & 75 \\ \hline
7 & Lowe's & LOW & Consumer Discretionary & 72 \\ \hline
8 & ADP & ADP & Information Technology & 70 \\ \hline
9 & UPS & UPS & Industrials & 67 \\ \hline
\end{tabular}
\label{tab:category2}
\end{table}

\subsubsection*{Mastercard}
\textbf{Ticker:} MA \\
\textbf{Sector:} Financials

Mastercard fits Category 2 because, by the end of 2022, AI was already embedded in fraud prevention, security, analytics, and personalization, but remained secondary to the broader payments-network business. AI supported transaction security, recommendation systems, and service differentiation, which gave it real commercial relevance. However, the firm’s revenues continued to depend overwhelmingly on payment volumes, network services, and transaction-related fees rather than on AI-specific demand. This means that AI had a visible and useful role in the company’s operations and product offerings, but it was not a major growth engine or a central determinant of the business model. Category 2 is therefore the most appropriate classification.

\textbf{Assessment:} Direct AI revenue exposure: None; Indirect AI demand exposure: Low; Strategic importance of AI: Medium; Operational reliance on AI: Medium; Economic materiality of AI: Low.

\textbf{Borderline note:} Mastercard belongs in Category 2 rather than Category 3 because AI supported commercial offerings and security functions, but was not economically meaningful enough to materially shape firm-wide demand or growth. It is above Category 1 because AI was clearly present in named products and services.

\textbf{Confidence score:} 84

\subsubsection*{Bank of America}
\textbf{Ticker:} BAC \\
\textbf{Sector:} Financials

Bank of America belongs in Category 2 because AI was already visible at scale in digital customer service by the end of 2022, especially through Erica, its AI-driven virtual financial assistant. This shows that AI had become commercially relevant in improving customer interaction and service delivery. Even so, the bank remained fundamentally a diversified financial institution whose revenues depended on lending, deposits, card activity, wealth management, and markets businesses rather than on AI-related products. AI was therefore important as a support technology for customer engagement and operational efficiency, but not central to the economics of the company. That makes Bank of America a clear example of low rather than moderate AI dependency.

\textbf{Assessment:} Direct AI revenue exposure: None; Indirect AI demand exposure: Low; Strategic importance of AI: Medium; Operational reliance on AI: Medium; Economic materiality of AI: Low.

\textbf{Borderline note:} Bank of America belongs in Category 2 rather than Category 3 because AI improved customer interaction and digital servicing, but it was not economically important enough to shape the bank’s overall demand or growth profile. It remains above Category 1 because AI had a real customer-facing role at scale.

\textbf{Confidence score:} 82

\subsubsection*{American Express}
\textbf{Ticker:} AXP \\
\textbf{Sector:} Financials

American Express fits Category 2 because AI and machine learning were already being used across servicing, analytics, transaction processing, and risk management by the end of 2022, but not as a major driver of the firm’s economic model. AI supported customer interactions, operational processes, and risk-related functions, making it strategically useful and commercially relevant. However, the company remained fundamentally a payments, card, and network-services business whose economics depended on spending volumes, merchant acceptance, cardmember loans, and fee income. This means AI functioned mainly as an enabling and efficiency-enhancing capability rather than as a material source of demand or growth. Category 2 is therefore the appropriate placement.

\textbf{Assessment:} Direct AI revenue exposure: None; Indirect AI demand exposure: Low; Strategic importance of AI: Medium; Operational reliance on AI: Medium; Economic materiality of AI: Low.

\textbf{Borderline note:} American Express belongs in Category 2 rather than Category 3 because AI was useful across several business functions, but not in a way that made it economically meaningful as a growth driver. It is above Category 1 because AI was already part of externally relevant servicing and risk functions.

\textbf{Confidence score:} 80

\subsubsection*{Starbucks}
\textbf{Ticker:} SBUX \\
\textbf{Sector:} Consumer Discretionary

Starbucks belongs in Category 2 because AI was already being used in personalization and selected operational tools by the end of 2022, but remained secondary to the firm’s retail coffee business. AI supported customer targeting, digital engagement, and store-level process improvements, including efforts to improve wait times and tailor the customer experience. These applications gave AI practical commercial value, especially in digital channels and loyalty-related activities. Nevertheless, Starbucks’ core economics still depended on beverage sales, store traffic, pricing, and brand strength rather than on AI-based monetization. AI therefore acted mainly as a supporting capability that improved execution and customer experience, which places the firm in Category 2 rather than a higher category.

\textbf{Assessment:} Direct AI revenue exposure: None; Indirect AI demand exposure: Low; Strategic importance of AI: Medium; Operational reliance on AI: Medium; Economic materiality of AI: Low.

\textbf{Borderline note:} Starbucks belongs in Category 2 rather than Category 3 because AI improved personalization and operating efficiency, but was not a meaningful independent source of demand or growth. It is above Category 1 because AI was already tied to customer-facing applications and store optimization.

\textbf{Confidence score:} 78

\subsubsection*{McDonald's}
\textbf{Ticker:} MCD \\
\textbf{Sector:} Consumer Discretionary

McDonald's fits Category 2 because pre-2023 public information showed visible AI-related activity in personalization and automated ordering, but these technologies remained peripheral to the core franchised restaurant model. AI was used mainly to improve customer experience, digital ordering, and selected operating processes. This indicates that AI had practical relevance to service delivery and efficiency, but not to the basic economics of the company. McDonald's still depended overwhelmingly on restaurant sales, franchise fees, and real-estate-related income rather than on AI-driven products or demand. As a result, AI was present and useful, but not economically central. This makes McDonald's a strong example of low AI dependency.

\textbf{Assessment:} Direct AI revenue exposure: None; Indirect AI demand exposure: Low; Strategic importance of AI: Medium; Operational reliance on AI: Low; Economic materiality of AI: Low.

\textbf{Borderline note:} McDonald's belongs in Category 2 rather than Category 3 because AI was mainly an operational and digital enhancement rather than a commercially meaningful business driver. It remains above Category 1 because the company had already deployed AI-related tools in customer-facing functions.

\textbf{Confidence score:} 77

\subsubsection*{CVS Health}
\textbf{Ticker:} CVS \\
\textbf{Sector:} Health Care

CVS Health belongs in Category 2 because AI was already being used to support insights, automation, and workflow efficiency by the end of 2022, but not as a central determinant of business performance. AI contributed to process improvements and decision support across customer service, pharmacy operations, and clinical workflows. This gave AI operational relevance and some commercial importance, especially where it improved service delivery and reduced manual work. However, CVS remained fundamentally a healthcare, pharmacy, and insurance-related company whose revenues depended on prescription dispensing, pharmacy services, and broader health-related activities rather than on AI-specific demand. AI therefore functioned as a secondary support layer rather than an economically meaningful growth engine.

\textbf{Assessment:} Direct AI revenue exposure: None; Indirect AI demand exposure: Low; Strategic importance of AI: Medium; Operational reliance on AI: Medium; Economic materiality of AI: Low.

\textbf{Borderline note:} CVS Health belongs in Category 2 rather than Category 3 because AI supported service delivery and internal process quality, but did not yet appear economically meaningful enough to drive firm-level demand or growth. It remains above Category 1 because AI was already linked to revenue-generating operating segments.

\textbf{Confidence score:} 75

\subsubsection*{Lowe's}
\textbf{Ticker:} LOW \\
\textbf{Sector:} Consumer Discretionary

Lowe's fits Category 2 because AI was visible by the end of 2022, but mainly in store operations, visualization, and associate support rather than in the core economic model. AI-related technologies supported retail execution, collaboration, and customer assistance, making them relevant to store processes and service quality. However, the company remained fundamentally a home-improvement retailer whose results depended on merchandise sales, housing-related demand, pricing, and store execution rather than on AI-based monetization. AI therefore had practical operational value, but it was not commercially important enough to shape the overall business in a material way. This places Lowe's clearly in the low AI dependency category.

\textbf{Assessment:} Direct AI revenue exposure: None; Indirect AI demand exposure: Low; Strategic importance of AI: Low; Operational reliance on AI: Medium; Economic materiality of AI: Low.

\textbf{Borderline note:} Lowe's belongs in Category 2 rather than Category 3 because AI was mainly a supporting operational technology rather than a commercially meaningful driver of demand or growth. It remains above Category 1 because AI had already moved beyond generic experimentation into identified retail applications.

\textbf{Confidence score:} 72

\subsubsection*{ADP}
\textbf{Ticker:} ADP \\
\textbf{Sector:} Information Technology

ADP belongs in Category 2 because AI and machine learning were already relevant enhancements to its data-driven human capital management solutions by the end of 2022, but not a dominant economic driver. AI supported insights, automation, and product value creation in payroll and HCM services, which made it strategically useful and commercially visible. However, ADP remained fundamentally a payroll, HCM, and employer-services company with revenues tied to recurring software, outsourcing, and business-process services rather than to AI-specific demand. This means AI improved the value proposition of existing offerings, but did not materially alter the core sources of revenue or demand. Category 2 is therefore the most defensible classification.

\textbf{Assessment:} Direct AI revenue exposure: None; Indirect AI demand exposure: Low; Strategic importance of AI: Medium; Operational reliance on AI: Medium; Economic materiality of AI: Low.

\textbf{Borderline note:} ADP belongs in Category 2 rather than Category 3 because AI was presented as a useful extension of broader HCM and data capabilities, not as an economically meaningful business driver in its own right. It remains above Category 1 because AI was explicitly linked to product value.

\textbf{Confidence score:} 70

\subsubsection*{UPS}
\textbf{Ticker:} UPS \\
\textbf{Sector:} Industrials

UPS is a lower-confidence but still defensible Category 2 firm. By the end of 2022, AI-related tools were already relevant to selected shipping, routing, and logistics support functions, especially where advanced algorithms improved real-time information and operating efficiency. This indicates that AI had practical value in optimizing parts of the logistics network and customer experience. Even so, UPS remained fundamentally a transportation and logistics company whose economics depended on shipment volumes, network density, pricing, labor, and fuel costs rather than on AI-based demand. AI therefore acted mainly as a secondary operational tool that improved execution, rather than as a commercially meaningful driver of growth. That places UPS appropriately in Category 2.

\textbf{Assessment:} Direct AI revenue exposure: None; Indirect AI demand exposure: Low; Strategic importance of AI: Low; Operational reliance on AI: Medium; Economic materiality of AI: Low.

\textbf{Borderline note:} UPS belongs in Category 2 rather than Category 3 because AI improved logistics execution and customer tools, but there is little evidence that it was economically meaningful enough to influence the overall growth profile of the business. It remains above Category 1 because AI had identified real-world operational applications by 2022.

\textbf{Confidence score:} 67

\subsubsection*{Conclusion on Category 2}

These firms can be used as anchor firms for the low AI-dependency category in the empirical analysis. Importantly, they were selected using only information available by 31 December 2022 and not on the basis of later stock price performance. The purpose of Category 2 is therefore to capture firms for which AI was already visible and useful before the generative-AI boom, but still clearly secondary to the core business model. This makes the category a useful comparison group between the moderate and high AI-dependency firms and the minimal AI-dependency firms.

\subsection{Category 1: Minimal AI Dependency}

Category 1 contains firms for which AI had little or no meaningful effect by 31 December 2022 on the company’s business model, revenues, demand, or near-term growth outlook. In practice, these are firms whose disclosed economics were driven overwhelmingly by regulated utility operations, physical infrastructure, property rents, or other non-AI commercial drivers. In this category, AI may still exist somewhere inside the organization, but it does not appear economically material to the firm.

The classification is intentionally conservative. Firms were included only when the available pre-2023 evidence suggested that AI was not economically meaningful to the business. If AI appeared as a visible secondary tool, feature, or efficiency enhancer, the firm was assigned to Category 2 instead. Applying this standard results in nine S\&P 500 firms being placed in Category 1: Southern Company, Duke Energy, Consolidated Edison, American Water Works, WEC Energy Group, Atmos Energy, Realty Income, Public Storage, and Simon Property Group.

\begin{table}[H]
\centering
\caption{Category 1 firms: Minimal AI dependency}
\begin{tabular}{|c|l|c|l|c|}
\hline
\textbf{Rank} & \textbf{Company} & \textbf{Ticker} & \textbf{Sector} & \textbf{Confidence score} \\ \hline
1 & Southern Company & SO & Utilities & 95 \\ \hline
2 & Duke Energy & DUK & Utilities & 93 \\ \hline
3 & Consolidated Edison & ED & Utilities & 91 \\ \hline
4 & American Water Works & AWK & Utilities & 90 \\ \hline
5 & WEC Energy Group & WEC & Utilities & 89 \\ \hline
6 & Atmos Energy & ATO & Utilities & 87 \\ \hline
7 & Realty Income & O & Real Estate & 85 \\ \hline
8 & Public Storage & PSA & Real Estate & 82 \\ \hline
9 & Simon Property Group & SPG & Real Estate & 80 \\ \hline
\end{tabular}
\label{tab:category1}
\end{table}

\subsubsection*{Southern Company}
\textbf{Ticker:} SO \\
\textbf{Sector:} Utilities

Southern Company is a very strong Category 1 firm because its disclosed 2022 economics were centered on regulated electricity and gas operations rather than on software or AI-linked demand. The company’s business model was fundamentally based on energy generation, transmission, distribution, and rate-based investment. This makes the firm’s economic logic overwhelmingly physical and regulatory rather than digital or model-driven. AI may have existed in operational support functions, but it did not appear to be commercially meaningful to revenues, demand, or near-term growth by the end of 2022. Southern Company therefore serves as a clean minimal-AI-dependency anchor firm in the sample.

\textbf{Assessment:} Direct AI revenue exposure: None; Indirect AI demand exposure: None; Strategic importance of AI: None; Operational reliance on AI: Low; Economic materiality of AI: None.

\textbf{Borderline note:} Southern Company belongs in Category 1 rather than Category 2 because its disclosed economics were overwhelmingly tied to regulated electricity and gas delivery. There is no clear evidence that AI had even secondary commercial significance by the end of 2022.

\textbf{Confidence score:} 95

\subsubsection*{Duke Energy}
\textbf{Ticker:} DUK \\
\textbf{Sector:} Utilities

Duke Energy fits Category 1 because its disclosed business model in 2022 was driven primarily by regulated utility operations, customer growth, rate-base expansion, and infrastructure investment rather than by technology-led product demand. The company’s commercial logic rested on electricity generation, distribution, and regulated returns across multiple utility jurisdictions. This is a classic physical infrastructure and regulation-based business model. Even if Duke Energy used modern digital systems internally, AI did not appear to have a meaningful effect on demand, revenues, or strategic positioning by the end of 2022. That makes Duke Energy a strong example of minimal AI dependency in the S\&P 500.

\textbf{Assessment:} Direct AI revenue exposure: None; Indirect AI demand exposure: None; Strategic importance of AI: None; Operational reliance on AI: Low; Economic materiality of AI: None.

\textbf{Borderline note:} Duke Energy belongs in Category 1 rather than Category 2 because the disclosed growth logic is regulatory and infrastructure-driven, not technology-driven. AI may exist in back-office processes, but it is not an economically meaningful feature of the firm’s business model.

\textbf{Confidence score:} 93

\subsubsection*{Consolidated Edison}
\textbf{Ticker:} ED \\
\textbf{Sector:} Utilities

Consolidated Edison belongs in Category 1 because its pre-2023 business model was clearly anchored in regulated electric, gas, and steam utility services. The company’s revenues and growth outlook depended primarily on rate cases, customer service obligations, utility infrastructure, and the provision of essential services rather than on any technology-driven commercial engine. This makes the firm’s economics heavily asset-based and regulated. AI may have had limited internal relevance for operating efficiency, but there is no indication that it materially influenced revenues, demand, or strategy by 31 December 2022. Consolidated Edison therefore fits comfortably in the minimal AI dependency category.

\textbf{Assessment:} Direct AI revenue exposure: None; Indirect AI demand exposure: None; Strategic importance of AI: None; Operational reliance on AI: Low; Economic materiality of AI: None.

\textbf{Borderline note:} Consolidated Edison belongs in Category 1 rather than Category 2 because its disclosed economics were almost entirely tied to regulated utility services and rate structures, not to AI-enabled products or efficiency layers with meaningful commercial importance.

\textbf{Confidence score:} 91

\subsubsection*{American Water Works}
\textbf{Ticker:} AWK \\
\textbf{Sector:} Utilities

American Water Works is a clear Category 1 firm because its business model was overwhelmingly defined by regulated water and wastewater services by the end of 2022. Its core economics depended on operating utility systems, serving regulated customer bases, and earning returns through traditional infrastructure and service provision. This is a network utility model with little connection to AI-driven demand or monetization. Any use of digital tools would have been supportive rather than economically defining. On the evidence available by 31 December 2022, AI did not have a meaningful effect on revenues, demand, or strategic positioning. American Water Works is therefore a strong anchor firm for minimal AI dependency.

\textbf{Assessment:} Direct AI revenue exposure: None; Indirect AI demand exposure: None; Strategic importance of AI: None; Operational reliance on AI: Low; Economic materiality of AI: None.

\textbf{Borderline note:} American Water Works belongs in Category 1 rather than Category 2 because its disclosed economics rest on regulated water and wastewater provision. There is no meaningful pre-2023 evidence that AI had even a secondary commercial role.

\textbf{Confidence score:} 90

\subsubsection*{WEC Energy Group}
\textbf{Ticker:} WEC \\
\textbf{Sector:} Utilities

WEC Energy Group fits Category 1 because its business model was clearly tied to electric and natural gas utility service rather than to AI-related products, demand, or monetization. The company’s economics were driven by regulated utility activities, infrastructure ownership, customer service, and the traditional performance drivers of a regional energy holding company. This makes WEC a highly physical and regulation-based firm. AI may have been present in small operational functions, but it did not appear commercially meaningful or strategically important by the end of 2022. For that reason, WEC Energy Group belongs in the minimal AI dependency category rather than in any higher category.

\textbf{Assessment:} Direct AI revenue exposure: None; Indirect AI demand exposure: None; Strategic importance of AI: None; Operational reliance on AI: Low; Economic materiality of AI: None.

\textbf{Borderline note:} WEC belongs in Category 1 rather than Category 2 because its disclosed business drivers are customer counts, utility regulation, and infrastructure ownership. AI does not appear to have had commercially meaningful relevance by end-2022.

\textbf{Confidence score:} 89

\subsubsection*{Atmos Energy}
\textbf{Ticker:} ATO \\
\textbf{Sector:} Utilities

Atmos Energy belongs in Category 1 because its 2022 business model was driven by natural gas distribution, pipeline, and storage assets rather than by software, data products, or AI-related demand. Its revenues and growth logic were tied to regulated customer service, system safety, reliability spending, and infrastructure investment. This is a highly physical and regulated business model with little reason to view AI as economically important. Even if some digital tools were used internally, AI did not appear to shape the company’s revenues, demand, or strategic outlook in a meaningful way by the end of 2022. Atmos Energy therefore fits well in the minimal AI dependency category.

\textbf{Assessment:} Direct AI revenue exposure: None; Indirect AI demand exposure: None; Strategic importance of AI: None; Operational reliance on AI: Low; Economic materiality of AI: None.

\textbf{Borderline note:} Atmos Energy belongs in Category 1 rather than Category 2 because the company’s disclosed performance drivers are pipeline safety, reliability, capital investment, and customer service in regulated gas markets. AI does not appear commercially meaningful.

\textbf{Confidence score:} 87

\subsubsection*{Realty Income}
\textbf{Ticker:} O \\
\textbf{Sector:} Real Estate

Realty Income is a strong Category 1 real-estate example because its business model is fundamentally based on long-term net lease agreements and rental cash flows from commercial properties. By the end of 2022, the company’s economics were clearly driven by property acquisition, tenant quality, lease duration, occupancy, and dividend-supporting real-estate income rather than by AI-related demand or software-based monetization. This makes the company’s commercial logic straightforwardly asset- and lease-based. AI may have supported back-office analysis or decision-making, but it did not appear economically meaningful to the firm’s revenues, tenant demand, or near-term growth profile. Realty Income therefore fits clearly in Category 1.

\textbf{Assessment:} Direct AI revenue exposure: None; Indirect AI demand exposure: None; Strategic importance of AI: None; Operational reliance on AI: Low; Economic materiality of AI: None.

\textbf{Borderline note:} Realty Income belongs in Category 1 rather than Category 2 because the firm’s revenue model is transparently lease- and property-based. AI may help internal decision-making, but it is not a disclosed commercial feature or economic driver.

\textbf{Confidence score:} 85

\subsubsection*{Public Storage}
\textbf{Ticker:} PSA \\
\textbf{Sector:} Real Estate

Public Storage fits Category 1 because its disclosed economics were overwhelmingly tied to self-storage real-estate operations by the end of 2022. The company’s performance depended on facility ownership, occupancy, rents, acquisitions, and same-store operating income rather than on AI-driven products or demand. This is a classic property-based business model in which the commercial engine is physical asset ownership and customer rental behavior. AI may have had some limited operational relevance, but it did not appear materially important to revenues, demand, or strategy. For that reason, Public Storage is a strong low-technology anchor firm and belongs in the minimal AI dependency category.

\textbf{Assessment:} Direct AI revenue exposure: None; Indirect AI demand exposure: None; Strategic importance of AI: None; Operational reliance on AI: Low; Economic materiality of AI: None.

\textbf{Borderline note:} Public Storage belongs in Category 1 rather than Category 2 because its disclosed business model is dominated by physical storage assets and rental income. There is no sign that AI had meaningful commercial relevance by the end of 2022.

\textbf{Confidence score:} 82

\subsubsection*{Simon Property Group}
\textbf{Ticker:} SPG \\
\textbf{Sector:} Real Estate

Simon Property Group belongs in Category 1 because its business model in 2022 was driven by traditional real-estate performance variables such as occupancy, property net operating income, redevelopment projects, and tenant demand across malls and outlet centers. The company’s value creation logic depended on physical assets, leasing, redevelopment, and portfolio management rather than on any AI-linked product, service, or demand channel. This makes the economic foundation of the firm clearly property-based. Even if digital tools were used internally, AI did not appear to have a meaningful effect on revenues, demand, or growth expectations by the end of 2022. Simon Property Group therefore fits appropriately in Category 1.

\textbf{Assessment:} Direct AI revenue exposure: None; Indirect AI demand exposure: None; Strategic importance of AI: None; Operational reliance on AI: Low; Economic materiality of AI: None.

\textbf{Borderline note:} Simon belongs in Category 1 rather than Category 2 because the disclosed economic story is almost entirely about occupancy, rental performance, redevelopment, and asset management. AI does not appear as even a secondary commercial driver.

\textbf{Confidence score:} 80

\subsubsection*{Conclusion on Category 1}

These firms can be used as anchor firms for the minimal AI-dependency category in the empirical analysis. Importantly, they were selected using only information available by 31 December 2022 and not on the basis of later stock price performance. The purpose of Category 1 is therefore to identify firms for which AI was not economically material before the generative-AI boom. This makes the category a useful baseline group when testing whether firms with greater pre-existing AI dependency experienced larger stock price increases during the later period.


\section{FORMER TEXT}

\section*{AI Dependency Classification Framework}


To examine whether firms with higher AI dependency experienced larger stock price increases during the generative-AI boom, the sample firms are classified into five categories based on the degree to which AI was economically material to the business. The purpose of this classification is to distinguish between firms for which AI was already central to revenues, demand, or business strategy before the boom, and firms for which AI played only a limited or peripheral role. The categorization is therefore intended to capture differences in firms’ \textit{ex ante} exposure to AI rather than differences in subsequent stock-market performance.

A key methodological principle is that the classification is based only on information that was publicly available by \textbf{31 December 2022}. This date is chosen to ensure that the classification predates the strongest phase of the generative-AI stock-market boom and is therefore not influenced by hindsight. The objective is to identify how exposed firms already were to AI before the market response intensified, rather than to classify firms based on what happened afterward. Consequently, the categorization does not rely on post-2022 stock returns, information released in 2023 or later, or media attention alone. Likewise, firms are not classified simply according to how often management mentioned ``AI,'' but rather according to whether AI appears to have had real economic significance for the company’s business model, revenues, demand, operations, or near-term growth outlook.

The classification is also deliberately conservative. Where the evidence is mixed, firms are assigned to the lower category. This reduces the risk of overstating AI exposure and helps preserve a clear distinction between firms for which AI was already economically central and firms for which it was still secondary. The resulting framework ranges from firms with very high AI dependency to firms with minimal AI dependency, thereby providing a structured basis for comparing stock-price developments across different levels of pre-existing AI exposure.

We have used AI to generate the five categories, from an in depth prompt where we have stressed the main attributed we wanted to catecterize by. The promt can be seen below: 

{\itshape

\subsection*{\textit{Prompt}}

You are helping with a master’s thesis in finance. The thesis studies whether firms with higher AI dependency experienced larger stock price increases during the generative-AI boom.

Your task is to choose around 10 S\&P 500 firms that best fit \textbf{AI Dependency Category [X]}, where \textbf{[X] = 1, 2, 3, 4, or 5}. You must choose the firms yourself and should not ask for company names.

\subsubsection*{\textit{Objective}}

Identify the firms that best belong in \textbf{Category [X]} using an \textit{ex ante} classification.

\subsubsection*{\textit{Methodological rule}}

Use only information that was publicly available by \textbf{31 December 2022}. The classification must reflect \textbf{economic dependency on AI}, not later market performance. Do not use stock returns after that date, information from 2023 or later, media hype alone, or simple counts of how often management said ``AI''.

\subsubsection*{\textit{AI dependency categories}}

\begin{description}
    \item[\textit{Category 5:}] Very high AI dependency. AI is central to the business model, demand, revenues, or near-term growth outlook.
    \item[\textit{Category 4:}] High AI dependency. AI is a major growth driver, but not the entire foundation of the business.
    \item[\textit{Category 3:}] Moderate AI dependency. AI is strategically relevant and economically meaningful, but not a dominant driver.
    \item[\textit{Category 2:}] Low AI dependency. AI is present mainly as a secondary tool, feature, or efficiency enhancer.
    \item[\textit{Category 1:}] Minimal AI dependency. AI has little or no meaningful effect on the company’s business model, revenues, demand, or operations.
\end{description}

\subsubsection*{\textit{Assessment criteria}}

Evaluate each firm based on five dimensions: direct AI revenue exposure, indirect exposure to AI-related demand, strategic importance of AI, operational reliance on AI, and overall economic materiality of AI.

\subsubsection*{\textit{Classification rules}}

Be conservative. If the evidence is mixed, choose the lower category. Do not assume that tech firms belong in high categories or that non-tech firms belong in low categories. Include only firms that clearly fit the category. It is acceptable to return fewer than 10 firms if fewer clearly qualify.

\subsubsection*{\textit{Required output}}

\begin{enumerate}
    \item \textbf{Category summary.} Start with the category number and label, a short explanation of what defines the category, and a short note on how strict the selection was.
    
    \item \textbf{Ranking table.} Provide a table with the following columns: Rank, Company, Ticker, Sector, and Confidence score.
    
    \item \textbf{Firm-by-firm analysis.} For each selected firm, provide the company name, ticker, and sector. Then write a short paragraph of 80--120 words explaining why the firm fits the category, using only information available by 31 December 2022. After that, report the following assessments: direct AI revenue exposure, indirect AI demand exposure, strategic importance of AI, operational reliance on AI, and economic materiality of AI, each using the scale \textit{High / Medium / Low / None}. Also include a brief borderline note explaining why the firm belongs in this category rather than the one above or below, and assign a confidence score from 0 to 100.
    
    \item \textbf{Final note.} End with a short paragraph explaining how the selected firms can be used as anchor firms for this category in a thesis dataset.
\end{enumerate}

\subsubsection*{\textit{Final instruction}}

Be precise, academically neutral, and consistent. Only include firms that genuinely fit \textbf{Category [X]} based on information available by 31 December 2022.

}

From this we get our companies split into our five categories:

\subsubsection*{Category 5 -- Very high AI dependency}

This category includes firms for which AI was already economically central by 31 December 2022. For these firms, AI was a core driver of the business model, revenues, demand, or near-term growth outlook. Their commercial performance depended directly on AI-related products, infrastructure, or monetization mechanisms rather than on broader activities in which AI played only a supporting role.

\begin{table}[H]
\centering
\caption{Category 5 firms: Very high AI dependency}
\small
\renewcommand{\arraystretch}{1.15}
\begin{tabularx}{\textwidth}{@{} C{1.1cm} L{4.2cm} C{1.6cm} L{4.3cm} C{2.3cm} @{}}
\toprule
\textbf{Rank} & \textbf{Company} & \textbf{Ticker} & \textbf{Sector} & \textbf{Confidence score} \\
\midrule
1 & NVIDIA & NVDA & Information Technology & 96 \\
2 & Alphabet & GOOGL & Communication Services & 88 \\
3 & Meta Platforms & META & Communication Services & 86 \\
\bottomrule
\end{tabularx}
\label{tab:category5}
\end{table}

\subsubsection*{Category 4 -- High AI dependency}

This category includes firms for which AI was already a major growth driver by the end of 2022, but not the entire foundation of the business. AI materially supported demand, product differentiation, commercialization, or customer adoption, yet the firm’s total revenue base still depended on broader platforms, ecosystems, or adjacent business lines. These firms therefore had substantial AI exposure, but were not as fully AI-centered as firms in Category 5.

\begin{table}[H]
\centering
\caption{Category 4 firms: High AI dependency}
\small
\renewcommand{\arraystretch}{1.15}
\begin{tabularx}{\textwidth}{@{} C{1.1cm} L{4.2cm} C{1.6cm} L{4.3cm} C{2.3cm} @{}}
\toprule
\textbf{Rank} & \textbf{Company} & \textbf{Ticker} & \textbf{Sector} & \textbf{Confidence score} \\
\midrule
1 & Microsoft & MSFT & Information Technology & 90 \\
2 & Amazon & AMZN & Consumer Discretionary & 88 \\
3 & AMD & AMD & Information Technology & 86 \\
4 & Adobe & ADBE & Information Technology & 84 \\
5 & Intuit & INTU & Information Technology & 81 \\
6 & ServiceNow & NOW & Information Technology & 79 \\
7 & Arista Networks & ANET & Information Technology & 74 \\
\bottomrule
\end{tabularx}
\label{tab:category4}
\end{table}

\subsubsection*{Category 3 -- Moderate AI dependency}

This category includes firms for which AI was strategically relevant and economically meaningful by 31 December 2022, but not a dominant driver of the business. In these firms, AI typically contributed through product enhancement, monetized features, workflow automation, or indirect support for demand. However, the firm’s revenues and growth outlook still depended primarily on broader software, services, semiconductor, or other core business activities rather than on AI itself.

\begin{table}[H]
\centering
\caption{Category 3 firms: Moderate AI dependency}
\small
\renewcommand{\arraystretch}{1.15}
\begin{tabularx}{\textwidth}{@{} C{1.1cm} L{4.2cm} C{1.6cm} L{4.3cm} C{2.3cm} @{}}
\toprule
\textbf{Rank} & \textbf{Company} & \textbf{Ticker} & \textbf{Sector} & \textbf{Confidence score} \\
\midrule
1 & Salesforce & CRM & Information Technology & 85 \\
2 & Oracle & ORCL & Information Technology & 83 \\
3 & Accenture & ACN & Information Technology & 81 \\
4 & Qualcomm & QCOM & Information Technology & 80 \\
5 & Palo Alto Networks & PANW & Information Technology & 79 \\
6 & Synopsys & SNPS & Information Technology & 78 \\
7 & Workday & WDAY & Information Technology & 76 \\
8 & Autodesk & ADSK & Information Technology & 74 \\
9 & Fortinet & FTNT & Information Technology & 73 \\
\bottomrule
\end{tabularx}
\label{tab:category3}
\end{table}

\subsubsection*{Category 2 -- Low AI dependency}

This category includes firms for which AI was clearly present by the end of 2022, but mainly as a secondary tool, feature, or efficiency enhancer. AI was used in areas such as personalization, fraud detection, customer service, logistics, or internal decision support, but it did not materially shape the firm’s core business model or revenue generation. AI therefore had practical relevance, but low overall economic materiality.

\begin{table}[H]
\centering
\caption{Category 2 firms: Low AI dependency}
\small
\renewcommand{\arraystretch}{1.15}
\begin{tabularx}{\textwidth}{@{} C{1.1cm} L{4.2cm} C{1.6cm} L{4.3cm} C{2.3cm} @{}}
\toprule
\textbf{Rank} & \textbf{Company} & \textbf{Ticker} & \textbf{Sector} & \textbf{Confidence score} \\
\midrule
1 & Mastercard & MA & Financials & 84 \\
2 & Bank of America & BAC & Financials & 82 \\
3 & American Express & AXP & Financials & 80 \\
4 & Starbucks & SBUX & Consumer Discretionary & 78 \\
5 & McDonald's & MCD & Consumer Discretionary & 77 \\
6 & CVS Health & CVS & Health Care & 75 \\
7 & Lowe's & LOW & Consumer Discretionary & 72 \\
8 & ADP & ADP & Information Technology & 70 \\
9 & UPS & UPS & Industrials & 67 \\
\bottomrule
\end{tabularx}
\label{tab:category2}
\end{table}

\subsubsection*{Category 1 -- Minimal AI dependency}

This category includes firms for which AI had little or no meaningful effect on the business model, revenues, demand, or operations by 31 December 2022. These firms were primarily driven by non-AI commercial factors such as regulated utility operations, physical infrastructure, property rents, or traditional consumer demand. Any use of AI was peripheral and not economically material to the firm’s overall performance.

\begin{table}[H]
\centering
\caption{Category 1 firms: Minimal AI dependency}
\small
\renewcommand{\arraystretch}{1.15}
\begin{tabularx}{\textwidth}{@{} C{1.1cm} L{4.2cm} C{1.6cm} L{4.3cm} C{2.3cm} @{}}
\toprule
\textbf{Rank} & \textbf{Company} & \textbf{Ticker} & \textbf{Sector} & \textbf{Confidence score} \\
\midrule
1 & Southern Company & SO & Utilities & 95 \\
2 & Duke Energy & DUK & Utilities & 93 \\
3 & Consolidated Edison & ED & Utilities & 91 \\
4 & American Water Works & AWK & Utilities & 90 \\
5 & WEC Energy Group & WEC & Utilities & 89 \\
6 & Atmos Energy & ATO & Utilities & 87 \\
7 & Realty Income & O & Real Estate & 85 \\
8 & Public Storage & PSA & Real Estate & 82 \\
9 & Simon Property Group & SPG & Real Estate & 80 \\
\bottomrule
\end{tabularx}
\label{tab:category1}
\end{table}

Taken together, these five categories provide an \textit{ex ante} measure of firm-level AI dependency that can be used to examine whether stock-price reactions during the generative-AI boom were systematically stronger among firms with greater pre-existing economic exposure to AI.

\end{document}
